% ========================================
% USAGE GUIDELINES
% ========================================
\section{Usage Guidelines for Researchers}

\subsection{Recommended Analysis Approaches}

\textbf{Cross-Country Comparisons:} Use the \texttt{broad\_fiscal\_gdp} variable to normalize for economic size differences. Be mindful of differences in fiscal accounting practices and reporting standards across countries.

\textbf{Temporal Analysis:} Combine \texttt{Year}, \texttt{Quarter}, and \texttt{Month} variables to capture timing and sequencing of policy responses. The highest policy concentration occurs in Q1--Q2 2020, reflecting initial pandemic response.

\textbf{Policy Type Analysis:} Utilize both \texttt{PolicyCode} and \texttt{PolicyCategory} to understand the mix of instruments employed. For robust analysis, we recommend using broad \texttt{PolicyCategory} classifications rather than fine-grained \texttt{PolicyCode} distinctions.

\textbf{Regional and Income Group Comparisons:} Leverage \texttt{region} and \texttt{income\_group} variables to identify patterns across economic development levels.

\subsection{Aggregation Considerations}

When calculating total fiscal commitments by country or period:

\begin{itemize}
\item Carefully consider how to handle observations with missing fiscal size values to avoid underestimation
\item Understand the \texttt{broad\_fiscal} classification before aggregating, as different values represent different fiscal impacts (direct spending vs. guarantees vs. deferrals)
\item Be aware that some policy measures may be recorded multiple times if announced in phases or modified over time
\item For guarantees and loans, recognize that reported amounts represent announced ceilings rather than actual utilization
\end{itemize}

\subsection{Potential Applications}

This dataset enables various research applications:

\begin{itemize}
\item Cross-country analysis of fiscal policy responses to the pandemic
\item Investigation of relationships between fiscal response timing, magnitude, and economic/health outcomes
\item Study of policy instrument choice across different institutional settings
\item Analysis of fiscal response heterogeneity by income level, region, or political economy factors
\item Examination of the composition of fiscal packages (direct spending vs. loans/guarantees vs. tax measures)
\end{itemize}

The dataset and documentation will be made publicly available upon publication. Researchers interested in early access should contact the author. The dataset is provided in Excel format (\texttt{.xlsx}) for maximum accessibility.

Please cite this dataset as follows:

\begin{quote}
[Author Name] (2025). ``COVID-19 Fiscal Response Database: A Comprehensive Micro-Level Dataset of Fiscal Policy Measures Across OECD Countries, 2020--2021.'' Working Paper, University of Basel.
\end{quote}